\documentstyle[%


righttag%


,11pt%


,amssymb%


,russian%               remove in Eng.


,izvrupar%                 Set izven% in Eng.


]{amsart}





%





\newcommand{\re}{\operatorname{Re}}


\newcommand{\tts}[1]{}





%\hoffset=-15mm%                  For Eng.


\setcounter{page}{20}


\god{2001}


\nomer{8~(471)} %                        Remove in Eng.


%\nomer{Vol.\,45, No.\,8}%               Set in Eng.


%\str{\pageref{begin}--\pageref{end}}%  Set in Eng.


\udk{517.546}





\author{\it Л.А.\,Аксентьев, В.П.\,Микка}


% NB:   ^^ remove in Eng.


\title{О поведении конформного радиуса\\ в подклассах


однолистных областей}





\date{}


%\pagestyle{myheadings}%         Set in Eng.


\pagestyle{plain} %              Remove in Eng.


\begin{document}


%\thispagestyle{plain}%          Set in Eng.


\maketitle


\label{begin}





В статье продолжаются исследования по поведению конформного


(внутреннего) радиуса односвязной области \cite{1}--\cite{4}. В


первом параграфе приводится формула для представления конформного


радиуса в малой окрестности фиксированной точки с выделением


представления для случая, когда точка является критической. Эта


формула позволяет дать новое доказательство известного условия


\cite{5} единственности критической точки в форме


$|\{f(\zeta),\zeta\}|\leq \frac{2}{(1-|\zeta|^2)^2}$


со шварцианом в левой части неравенства.


Во вторoм параграфе выписаны два семейства спиралеобразных функций с


нарушением единственности критической точки соответствующего


конформного радиуса и сформулированы две теоремы о единственности


критической точки конформного радиуса для подклассов


спиралеобразных областей.


В третьем параграфе содержится утверждение, дополняющее результаты из


\cite{4}, связанные с поведением конформного радиуса


почти выпуклых областей.





\section{}





Конформный радиус области $D$ в точке


$z$ вычисляется с помощью функции $z=f(\zeta)$, конформно


отображающей круг $E=\{\zeta:|\zeta|<1 \}$ на область $D$,


$$


R({D},z)=R(f(E),f(\zeta))=|f'(\zeta)|(1-|\zeta|^2)\equiv R(\zeta).


$$


Получим разложение этой вещественной положительной (т.\,к.


$f'(\zeta)\neq 0$, $\zeta\in E$) функции в окрестности точки $a$,


переписав последнее равенство из предыдущей формулы в виде


\begin{equation}\label{1}


R(a+\rho e^{i\theta})=|f'(a+\rho e^{i\theta})|(1-|a+\rho


e^{i\theta}|^2).


\end{equation}


Для второго множителя в (\ref{1}) имеем


\begin{equation}\label{2}


1-(a+\rho e^{i\theta})(\overline{a}+\rho e^{-i\theta})=


1-|a|^2-2\re(\overline{a}


e^{i\theta})\rho-\rho^2.


\end{equation}





Первый множитель запишем в виде $|f'(a+\rho


e^{i\theta})|=\exp(\ln|f'(a+\rho e^{i\theta})|)$ и разложим вначале функцию


$$


\ln f'(a+\rho e^{i\theta})=\ln


f'(a)+\frac{f''(a)}{f'(a)}e^{i\theta}\rho


+\frac{1}{2} \bigg( \frac{f'''(a)}{f'(a)}-\bigg(\frac{{f''}(a)}{{f'}(a)}\bigg)^2 \bigg)


e^{i2\theta}\rho^2+O(\rho^3).


$$


Это разложение приводит к представлению


$$


\ln |f'(a+\rho e^{i\theta})|= \re \ln f'(a+\rho e^{i\theta})


=\ln |f'(a)|+a_1\rho+a_2\rho^2+O(\rho^3),


$$


из которого следует


$$


|f'(a+\rho


e^{i\theta})|=|f'(a)|e^{a_1\rho+a_2\rho^2+O(\rho^3)}


=|f'(a)|\bigg(1+a_1\rho+\bigg(a_2+\frac{a_1^2}{2}\bigg)\rho^2+O(\rho^3)\bigg),


$$


где


$$


a_1=\re \bigg(\frac{f''(a)}{f'(a)}e^{i\theta}\bigg),\quad


a_2=\frac{1}{2}\re\bigg[ \bigg(


\frac{f'''(a)}{f'(a)}-\bigg(\frac{{f''}(a)}{{f'}(a)}\bigg)^2


\bigg)e^{i2\theta}\bigg].


$$


Поэтому с учетом (\ref{2}) будем иметь


\begin{equation}\label{3}


R(a+\rho e^{i\theta})=R(a)\:(1+b_1\rho+b_2\rho^2+O(\rho^3)),


\end{equation}


причем


\begin{gather}\label{4}


b_1=\re B_1,\quad B_1=\bigg( \frac{f''(a)}{f'(a)}-


\frac{2\overline{a}}{1-|a|^2} \bigg) e^{i\theta},\\


%


b_2=\frac{1}{2}\re\bigg(\bigg[


\frac{f'''(a)}{f'(a)}-\bigg(\frac{{f''}(a)}{{f'}(a)}\bigg)^2


\bigg]e^{i2\theta}\bigg)


+\frac{1}{2} \bigg[ \re\bigg(\frac{f''(a)}{f'(a)}e^{i\theta}\bigg)


\bigg]^2 - \notag \\


%


-\re\bigg(\frac{f''(a)}{f'(a)}e^{i\theta}\bigg)


\frac{2\re{(\overline{a}e^{i\theta})}}{1-|a|^2} -\frac{1}{1-|a|^2}.\notag


\end{gather}





Представим коэффициент $b_2$ в форме


\begin{equation}\label{5}


b_2=\frac12 \re\bigg(\{f(a),a


\}e^{i2\theta}-\frac{2}{(1-|a|^2)^2}\bigg)+c_2,


\end{equation}


в которую входят шварциан


$\{f(a),a\}=\frac{f'''(a)}{f'(a)}-\frac32\Big(\frac{f''(a)}{f'(a)}


\Big)^2$ и коэффициент


$$


c_2=\frac14\re\bigg(\frac{f''(a)}{f'(a)}e^{i\theta}


\bigg)^2+\frac12\bigg[\re\bigg(\frac{f''(a)}{f'(a)}e^{i\theta}\bigg)\bigg]^2-


\re\bigg(\frac{f''(a)}{f'(a)}e^{i\theta}


\bigg)\frac{2\re(\overline{a}e^{i\theta})}{1-|a|^2}+\frac{|a|^2}{(1-|a|^2)^2}.


$$





Преобразуем громоздкое выражение для этого коэффициента, записав


второе равенство из (\ref{4}) в виде


$$


a\bigg(\frac{f''(a)}{f'(a)}-\frac{2\overline{a}}{1-|a|^2}\bigg)e^{i\theta}=a


B_1\ \Rightarrow \


a\frac{f''(a)}{f'(a)}=aB_1e^{-i\theta}+\frac{2|a|^2}{1-|a|^2}.


$$


Отсюда следует равенство квадратов модулей


$$


|a|^2\bigg|\frac{f''(a)}{f'(a)}\bigg|^2=|a|^2|B_1|^2+


\frac{4|a|^2}{1-|a|^2}\re(a B_1


e^{-i\theta})+\frac{4|a|^4}{(1-|a|^2)^2}


$$


и после сокращения на $|a|^2$


\begin{equation}\label{6}


\bigg|\frac{f''(a)}{f'(a)}\bigg|^2=|B_1|^2+\frac{4\re(a B_1


e^{-i\theta})}{1-|a|^2}+\frac{4|a|^2}{(1-|a|^2)^2}.


\end{equation}


Учтем равенства (\ref{4}) и (\ref{6}), а также соотношение


$$


\re(u+iv)^2-2u^2=-|u+iv|^2


$$


для вещественных величин $u$, $v$, причем


$$


u+iv=e^{i\theta}f''(a)/f'(a),


$$


в несложных преобразованиях коэффициента $c_2$. Последовательно


будем иметь


\begin{gather*}


c_2=\frac14\re \bigg(\frac{f''(a)}{f'(a)}e^{i\theta}\bigg)^2-


\frac12\bigg[\re\bigg(\frac{f''(a)}{f'(a)}e^{i\theta}\bigg)\bigg]^2+\\


+\re \bigg(\frac{f''(a)}{f'(a)}e^{i\theta}\bigg)


\bigg[\re \bigg(\frac{f''(a)}{f'(a)}e^{i\theta}\bigg)-


\frac{2\re(\overline{a}e^{i\theta})}{1-|a|^2}\bigg]+


\frac{|a|^2}{(1-|a|^2)^2}=\\


=-\frac{1}{4}\bigg|\frac{f''(a)}{f'(a)}e^{i\theta}\bigg|^2+


b_1\re \bigg(\frac{f''(a)}{f'(a)}e^{i\theta}\bigg)+


\frac{|a|^2}{(1-|a|^2)^2}


\overset{\text{(см. (\ref{6}))}}{=}\\=


-\frac{|B_1|^2}{4}-\frac{\re(aB_1e^{-i\theta})}{1-|a|^2}-


\frac{|a|^2}{(1-|a|^2)^2} +


b_1\re \bigg(e^{i\theta}\frac{f''(a)}{f'(a)}\bigg)


+\frac{|a|^2}{(1-|a|^2)^2}\Rightarrow


\end{gather*}


\begin{equation}\label{7}


c_2=-\frac{|B_1|^2}{4}-\frac{\re(aB_1e^{-i\theta})}{1-|a|^2}+


b_1\re \bigg(e^{i\theta}\frac{f''(a)}{f'(a)}\bigg).


\end{equation}





Конформный радиус наглядно характеризуется гладкой поверхностью с


уравнением $\Omega=R(\zeta)$ над горизонтальным кругом $E$ в


трехмерном пространстве с вертикальной осью $\Omega$. Точка из


круга $E$, над которой лежит точка поверхности конформного радиуса


с касательной плоскостью, параллельной горизонтальной плоскости,


называется критической точкой. Поверхность конформного радиуса


$\widetilde{\Omega}=R(f(E),z)$ можно построить и над областью


${D}=f(E)$. Хотя эти поверхности и будут различными, но


количество критических точек на них будет одинаковым.





Докажем, что любой критической точке $\zeta_0$ конформного радиуса


$R(f(E),f(\zeta))$ над кругом $E$ ставится во взаимно однозначное


соответствие критическая точка $z_0=f(\zeta_0)$ конформного


радиуса над областью ${D}=f(E)$.


Действительно, уравнения для определения этих критических точек


$\frac{\partial R}{\partial \zeta}=0$ и $\frac{\partial


R}{\partial z}=0$ можно переписать следующим образом:


$$


\frac{\partial R}{\partial \zeta}=0\ \Rightarrow


\ \frac{\partial R}{\partial z}=


\frac{\partial R}{\partial \zeta}\frac{\partial \zeta}{\partial z}+


\frac{\partial R}{\partial \overline{\zeta}}\frac{\partial \overline{\zeta}}{\partial


z}=0.


$$


Но $\frac{\partial \overline{\zeta}}{\partial z}=0$ в силу того,


что $\zeta=f^{-1}(z)$ --- аналитическая функция по $z$, а


$\overline{\zeta}=\overline{f^{-1}(z)}$ --- аналитическая функция по


$\overline{z}$, поэтому выполняется комплексное условие Коши--Римана.


Значит, уравнения $\frac{\partial R}{\partial \zeta}=0$ и $\frac{\partial


R}{\partial z}= \frac{\partial R}{\partial \zeta}\frac{\partial


\zeta}{\partial z}=0$ отличаются на функциональный множитель,


отличный от нуля для локально однолистных отображений $f(\zeta)$ и


$f^{-1}(z)$.





Необходимым и достаточным условием наличия критической точки


является равенство нулю коэффициента (\ref{4}) в разложении


(\ref{3}). Это разложение в окрестности критической точки примет


вид


\begin{equation}\label{8}


R(a+\rho e^{i\theta})=R(a)(1+\widehat{b}_2(\theta)\rho^2+O(\rho^3)),


\end{equation}


причем $b_1=0$, $0\leq \theta <2\pi$, $\Leftrightarrow$


$B_1=0\Rightarrow c_2=0$ (в силу (\ref{7})) и поэтому


\begin{equation}\label{9}


\widehat{b}_2(\theta)=\frac12\re\bigg(\{f(a),a\} e^{i2\theta}-


\frac{2}{(1-|a|^2)^2}\bigg).


\end{equation}





Частный вид разложений (\ref{3}), (\ref{8}) при $a=0$ и $f'(0)=1$


был получен и применен в \cite{2}, \cite{3}. Именно, справедливо


равенство, несколько измененное по сравнению с \cite{2}, \cite{3},


\begin{gather*}


R(\rho e^{i\theta})=1+\re(f''(0)e^{i\theta})\rho+


\bigg\{\frac12\re[(f'''(0)-f''{}^2(0))e^{i2\theta}]+\\


+\frac12[\re(f''(0)e^{i\theta})]^2-1\bigg\}\rho^2+O(\rho^3)=\\


=1+2\re(a_2e^{i\theta})\rho+\{\re[(3a_3-2a_2^2)e^{i2\theta}]+


2[\re(a_2e^{i\theta})]^2-1\}\rho^2+O(\rho^3).


\end{gather*}





С помощью представления (\ref{8}), (\ref{9}) можно выяснить


структуру поверхности конформного радиуса над окрестностью


критической точки $a$. Соотношение $b_1=0$ при всех $\theta$,


$0\leq\theta <2\pi$, из (\ref{4}) приводит к уравнению


\begin{equation}\label{10}


\frac{f''(a)}{f'(a)}=\frac{2\overline{a}}{1-|a|^2},


\end{equation}


которому должна удовлетворять критическая точка $a$.





Если $|\{f(a),a\}|\leq \frac{2}{(1-|a|^2)^2}$, то


$\widehat{b}_2(\theta)\leq 0$, поэтому $a$ будет точкой максимума


поверхности $\Omega=R(\zeta)$ или полуседловой точкой (с


изменением знака кривизны хотя бы для одной линии в срезax


поверхности). При $|\{f(a),a\}|> \frac{2}{(1-|a|^2)^2}$


коэффициент $\widehat{b}_2(\theta)$ будет иметь переменный знак,


поэтому $a$ не будет точкой максимума. Эти эффекты приводят к двум


важным выводам.


\medskip





$1^0$. Если выполняется условие однолистности Нехари \cite{6}


\begin{equation}\label{11}


|\{f(\zeta),\zeta\}|\leq \frac{2}{(1-|\zeta|^2)^2},\quad \zeta\in E,


\end{equation}


то у поверхности конформного радиуса будет одна критическая точка


--- единственный максимум \cite{5}, за исключением случая, когда


$f(E)$ является полосой.





Представим новое (совсем короткое) доказательство этого факта.


Действительно, предположение о существовании двух критических


точек приводит к противоречию. Для получения противоречия нужно


обе критические точки расположить на вещественном диаметре с


помощью дополнительного дробно-линейного преобразования, не


нарушающего условия (\ref{11}). Эти критические точки являются


точками максимума поверхности $R(\zeta)$ в силу (\ref{11}). При


строгом неравенстве (\ref{11}) такое невозможно, т.\,к. между


двумя максимумами должен быть хотя бы один минимум на вертикальном


срезе поверхности $\Omega=R(\zeta)$, а это не допускает строгое


неравенство (\ref{11}). Если же выполняется нестрогое неравенство


(\ref{11}), то появление двух максимумов над действительным


диаметром приведет к сплошному максимуму между двумя отмеченными


максимумами, т.\,е. к континууму критических точек, которые будут


удовлетворять уравнению


$$


\frac{f''(t)}{f'(t)}=\frac{2{t}}{1-t^2}.


$$


При нормировке $f'(0)=1$, $f(0)=0$ этому уравнению удовлетворяет


единственная функция $f(t)=\frac12\ln\frac{1+t}{1-t}$. Она


определена во всем круге $E$ и отображает его на горизонтальную


полосу шириной $\frac{\pi}{2}$. Опровержение возможности точки


перегиба вместо второго максимума можно дать при разложении


$1/R({D},z)$.


\medskip





$2^0$. Как было отмечено, при выполнении условия


\begin{equation}\label{12}


|\{f(a),a\}|> \frac{2}{(1-|a|^2)^2}


\end{equation}


точка $a$ не будет точкой максимума. Так как при естественных


ограничениях (напр., при непрерывности $f(\zeta)$ в


$\overline{E}$ \cite{1}) хотя бы один максимум у $R(\zeta)$ всегда


существует, то условие (\ref{12}) в критической точке $a$ является


гарантией существования более одной критической точки поверхности


конформного радиуса. Поэтому выполнение условия (\ref{12}) для


конкретных функций означает, что конформные радиусы, связанные с


этими функциями, имеют не менее двух критических точек. Другое


обоснование неединственности критической точки при условии


(\ref{12}) дано в (\cite{7}, теорема 3) с использованием индексов


критических точек.





\section{}





По отношению к классу спиралеобразных функций


в круге $E$ с условием


\begin{equation}\label{13}


\re\bigg(e^{i\gamma}\zeta\frac{f'(\zeta)}{f(\zeta)}\bigg)\geq0,\quad


-\frac{\pi}{2}<\gamma<\frac{\pi}{2},


\end{equation}


два пробных семейства можно построить с использованием функции


$\big(\zeta=\frac{w-1}{w+e^{-i2\gamma}}\big)$


$$


w=\frac{1+e^{-i2\gamma}\zeta}{1-\zeta}=


\frac{1+e^{-i\gamma}\zeta_1}{1-e^{i\gamma}\zeta_1},


$$


отображающей круг $E$ на полуплоскость $\re(we^{i\gamma})>0$


(которой принадлежит $1=w(0)$) с граничной прямой, проходящей


через точки $0$ и $i\sin\gamma e^{-i\gamma}$, т.\,к.


\begin{gather*}


\zeta=1\mapsto w=\infty,\ \;


\zeta=-e^{i2\gamma}\mapsto w=0,\\


\zeta=-1\mapsto w=(1-e^{-i2\gamma})/2=e^{-i\gamma}i\sin\gamma,


\end{gather*}


а для любой граничной точки на единичной окружности имеем


$$


\zeta=e^{i\theta}\mapsto w=\frac{e^{i(-\gamma+\theta/2)}


2\cos(\gamma-\frac{\theta}{2})}


{e^{i\theta/2}(-2i\sin\frac{\theta}{2})} =


e^{i(\pi/2-\gamma)}\frac{\cos(\gamma-\frac{\theta}{2})}


{\sin\frac{\theta}{2}}.


$$





Первое семейство было определено в \cite{4} с помощью равенства


\begin{equation}\label{14}


\zeta\frac{f'(\zeta)}{f(\zeta)}=


\frac{1+e^{-i\gamma}\beta\zeta^2}{1-e^{i\gamma}\alpha\zeta^2},


\quad \alpha,\beta\in(-1,1].


\end{equation}


Для этого семейства получен такой результат в \cite{4}. Шварциан при


$\zeta=0$ определяет соотношение


$$


|\{f(\zeta),\zeta \}| \big|_{\zeta=0}=


3|\alpha e^{i\gamma}+\beta e^{-i\gamma}|>2


\ \Rightarrow \


\alpha^2+\beta^2+2\alpha\beta\cos2\gamma>4/9.


$$


В плоскости $\alpha+i\beta$ получим множество ${\cal


D}(\gamma)$, которое является внешностью эллипса с


$\gamma\neq\pm{\pi}/{4}$, внешностью окружности (при


$\gamma=\pm{\pi}/{4}$), внешностью полосы с двумя параллельными


граничными прямыми (при $\gamma=0$ и $\gamma=\pm{\pi}/{2}$).


Принадлежность $\alpha+i\beta\in{\cal D}(\gamma)$ означает,


что конформный радиус для области $f(E)$ с $f(\zeta)$,


определяемой из уравнения (\ref{14}), будет иметь больше одной


критической точки.





Второе семейство пробных функций получится при такой


последовательной конструкции


\begin{multline*}


w=\bigg[\cos\gamma\frac{1+\beta\zeta}


{1-\alpha\zeta}+i\sin\gamma\bigg]e^{-i\gamma} =


\frac{\cos\gamma+i\sin\gamma+


(\beta\cos\gamma-i\alpha\sin\gamma)\zeta}{1-\alpha\zeta}e^{-i\gamma}=\\


=\frac{1+(\beta\cos\gamma-i\alpha\sin\gamma)


e^{-i\gamma}\zeta}{1-\alpha\zeta},


\end{multline*}


т.\,е.


\begin{equation}\label{15}


w=\frac{1+c\zeta}{1-\alpha\zeta},\quad


c=(\beta\cos\gamma-i\alpha\sin\gamma)e^{-i\gamma};\ \;


\alpha,\beta \in(-1,1],


\end{equation}


которая при $\alpha=\beta=1$ переходит в выражение


\begin{equation}\label{16}


w=\bigg(\cos\gamma


\frac{1+\zeta}{1-\zeta}+i\sin\gamma\bigg)e^{-i\gamma}=


\frac{\cos\gamma+i\sin\gamma+(\cos\gamma-i\sin\gamma)\zeta}


{1-\zeta}e^{-i\gamma}= \frac{1+e^{-i2\gamma}\zeta}{1-\zeta}.


\end{equation}


Геометрический смысл функции (\ref{16}) был выяснен выше.





Теперь с использованием (\ref{15}) функции второго семейства


определяются из уравнения


\begin{equation}\label{17}


\frac{\zeta


f'(\zeta)}{f(\zeta)}=\frac{1+c\zeta^2}{1-\alpha\zeta^2}, \quad


c=(\beta\cos\gamma-i\alpha\sin\gamma)e^{-i\gamma}.


\end{equation}


Отсюда легко получить


$$


\{f(\zeta),\zeta \}\big|_{\zeta=0}=


\bigg[\bigg(\frac{f''}{f'}


\bigg)'-\frac12\bigg(\frac{f''}{f'}


\bigg)^2\bigg]\bigg|_{\zeta=0}=3(\alpha+c).


$$


Оценка $|\{f(\zeta),\zeta\}\big|_{\zeta=0}|>2$ приводит к


неравенству $|\alpha+\beta|>{2}/(3\cos\gamma)$, т.\,к.


$c+\alpha=e^{-i\gamma}(\beta\cos\gamma+\alpha\cos\gamma)$. Здесь


аналогом множества ${\cal D}(\gamma)$ (для первого семейства


функций) будет внешность полосы с граничными прямыми


\begin{equation}\label{18}


\alpha+\beta=\pm{2}/(3\cos\gamma).


\end{equation}


При $\gamma=0$ получаются соответствующие прямые для


звездообразных функций \cite{4}.





По аналогии со звездообразными функциями \cite{4} имеются


основания предполагать, что внутри полосы с граничными линиями


(\ref{18}) найдутся области в плоскости параметров


$\alpha+i\beta$, гарантирующие единственность критической точки


конформного радиуса не только для функций с уравнением (\ref{17}),


но и для всех функций, определяемых подчинением


\begin{equation}\label{19}


\zeta\frac{f'(\zeta)}{f(\zeta)}\prec\frac{1+c\zeta}{1-\alpha\zeta},\ \;


|c+\alpha|=|e^{-i\gamma}(\alpha+\beta)\cos\gamma|\leq\frac23,


\end{equation}


причем постоянная $c$ представлена формулой (\ref{17}).





Займемся оценками, в результате которых получим одно из основных


утверждений статьи. Из условия подчинения (\ref{19}) следует


\begin{equation}\label{20}


\zeta\frac{f'(\zeta)}{f(\zeta)}=


\frac{1+c\varphi(\zeta)}{1-\alpha\varphi(\zeta)},\quad


\varphi(\zeta)\sim\zeta^2 \text{\; около \;} \zeta=0.


\end{equation}


С помощью логарифмической производной из соотношения (\ref{20})


имеем


\begin{gather*}


\frac{1}{\zeta}+\frac{f''}{f'}-\frac{f'}{f}=


\frac{(c+\alpha)\varphi'}{(1+c\varphi)(1-\alpha\varphi)}


\ \Rightarrow\\


\zeta\frac{f''}{f'}=(c+\alpha)


\bigg(\frac{\zeta\varphi'}{(1+c\varphi)(1-\alpha\varphi)}+


\frac{\varphi}{1-\alpha\varphi}\bigg).


\end{gather*}


Нетрудные оценки с использованием леммы Шварца (\cite{8}, гл.\,8,


\S\,1) приводят к неулучшаемому итоговому неравенству


\begin{gather*}


\bigg|\zeta\frac{f''(\zeta)}{f'(\zeta)}\bigg|\leq|c+\alpha|


\bigg(\frac{|\zeta|2|\zeta|(1-|\varphi|^2)(1-|\zeta|^4)^{-1}}


{|1+c\varphi||1-\alpha\varphi|}+


\frac{|\zeta|^2}{|1-\alpha\varphi|} \bigg)\leq\\


%


\leq |c+\alpha|\frac{|\zeta|^2}{1-|\zeta|^2}


\bigg(\frac{2(1+|\varphi|)}{1+|\zeta|^2}


\frac{1-|\varphi|}{|1+c\varphi|\,|1-\alpha\varphi|}+


\frac{1-|\zeta|^2}{1-|\varphi|}\:\frac{1-|\varphi|}{|1-\alpha\varphi|} \bigg)


\leq\\


%


\leq\frac{|\zeta|^2}{1-|\zeta|^2}|c+\alpha|


\bigg(2\frac{1-|\varphi|}{|1+c\varphi||1-\alpha\varphi|}+


\frac{1-|\varphi|}{|1-\alpha\varphi|}\bigg)


\overset{\text{def}}{=}\\


=\frac{|\zeta|^2}{1-|\zeta|^2}|c+\alpha|\Phi(\varphi)\leq


\frac{3|c+\alpha||\zeta|^2}{1-|\zeta|^2}.


\end{gather*}


В последнем неравенстве учтена оценка


\begin{gather*}


\Phi(\varphi)\leq


\frac{2(1-|\varphi|)}{(1-|c||\varphi|)(1-|\alpha|\,|\varphi|)}+


\frac{1-|\varphi|}{1-|\alpha||\varphi|}


\overset{\text{def}}{=}\\


=\Phi_1(|\varphi|)+\Phi_2(|\varphi|)\leq\Phi_1(0)+\Phi_2(0)=\Phi(0),


\end{gather*}


для обоснования которой покажем, что функции $\Phi_1(t)$ и


$\Phi_2(t)$ не возрастают на отрезке $[0,1]$ при условии


\begin{equation}\label{*}


|\alpha|+|c|\leq 1.


\end{equation}


Достаточно подсчитать и оценить производные этих функций. Получим


$$


\Phi_1'(t)=2\frac{|\alpha|\,|c|\big[(t-1)^2+(|\alpha|+


|c|-1-|\alpha|\,|c|)(|\alpha|\,|c|)^{-1}\big]}


{(1-|\alpha|t)^2(1-|c|t)^2}\leq 0,


$$


т.к. квадратная скобка в числителе оценивается при условии


(\ref{*}) величиной


$$


[(t-1)^2+(|\alpha|+|c|-1-|\alpha|\,|c|)(|\alpha|\,|c|)^{-1}]


\leq (t-1)^2-1\leq0,\quad t\in[0,1].


$$


Для второй функции подсчет совсем простой


$\Phi_2'(t)=\frac{-1+|\alpha|}{(1-|\alpha|t)^2}\leq0$


при $|\alpha|\leq1$.





Таким образом, справедлива оценка


$$


\bigg|\zeta


\frac{f''(\zeta)}{f'(\zeta)}\bigg|<\frac{2|\zeta|^2}{1-|\zeta|^2},\quad


\zeta\in E,\ \; \zeta \neq 0,


$$


если $|c+\alpha|<\frac23$ и $|c|+|\alpha|\leq 1$, что эквивалентно


неравенствам


\begin{equation}\label{21}


|\alpha+\beta|\cos\gamma<2/3 \text{\; и \;}


|\alpha|+\sqrt{\alpha^2\sin^2\gamma+\beta^2\cos^2\gamma}\leq 1.


\end{equation}


Поэтому уравнение (\ref{10}) не будет иметь решений, кроме $a=0$.





Тем самым доказана





\begin{thm}


Любая спиралеобразная функция $f(\zeta)$,


удовлетворяющая соотношению $(\ref{20})$, в котором параметры


подчинены условию $(\ref{21})$, обеспечивает единственность


критической точки в $0$ конформного радиуса области $f(E)$.


Увеличить правую часть в первом из неравенств $(\ref{21})$ нельзя в


силу $(\ref{17})$, $(\ref{18})$.


\end{thm}





На международной конференции 1999 года в Румынии был доложен


результат по спиралеобразным функциям, отраженный в тезисах


\cite{9}. Используем одну импликацию из этих тезисов:


$$


\re(e^{i\gamma}\zeta\frac{f'(\zeta)}{f(\zeta)})>0,\ \;


|\gamma|>{\pi}/{3},


\ \Rightarrow 


\ \bigg|\frac{f''(\zeta)}{f'(\zeta)}\bigg|<


2\frac{1+(4/3)\sin(2|\gamma|)}{1-|\zeta|^2},\quad \zeta\in E.


$$


В применении к функции ${\cal F}(\zeta)=f(\alpha\zeta)$ правое


неравенство перепишется


$$


\bigg|\frac{{\cal F}''(\zeta)}{{\cal F}'(\zeta)}\bigg|<


2\alpha\frac{1+(4/3)\sin(2|\gamma|)}{1-|\alpha\zeta|^2}<


\frac{2}{1-|\zeta|^2}


$$


в предположении, что


\begin{equation}\label{22}


\alpha<\big[1+(4/3)\sin(2|\gamma|)\big]^{-1}.


\end{equation}


Поэтому справедлива





\begin{thm} Если функция $f(\zeta)$ является


спиралеобразной в $E$, т.\,е. удовлетворяет неравенству


$\re\big(e^{i\gamma}\zeta\frac{f'(\zeta)}{f(\zeta)}\big)>0$, с


дополнительными условиями $\pi/3<|\gamma|< {\pi}/{2}$ и $f''(0)=0$, то


область $f(\alpha E)$ будет иметь конформный радиус с единственной


критической точкой в $0$ при выполнении неравенства $(\ref{22})$.


\end{thm}





\section{}





Для формулировки утверждения о конформном радиусе в подклассе


почти выпуклых функций понадобится теорема 1 из \cite{4}, которая


будет представлена как


\medskip





{\bf Лемма} \cite{4}. {\it


Области ${\cal D}_m$ в плоскости


$\alpha+i\beta$, где выполняется импликация


\begin{equation}\label{23}


\bigg\{f'(\zeta)\prec T(\alpha,\beta;\zeta)=


\frac{1+\beta\zeta}{1-\alpha\zeta},\ \; \alpha+i\beta \in {\cal D}_m\bigg\}


\ \Rightarrow 


\ \bigg|\frac{f''(\zeta)}{f'(\zeta)}\bigg|


\leq\frac{2m|\zeta|}{1-|\zeta|^2},\quad


0\leq m\leq 2,


\end{equation}


представляют собой трапеции с вершинами $1-i$, $1+i(m-1)$,


$m-1+i$, $-1+i$, когда $1\leq m<2$ $($при $m=2$ получается


треугольник$)$, и криволинейные трапеции с прямолинейными основаниями


$[\frac{m-1}{2}+i\frac{m+1}{2},\frac{m+1}{2}+i\frac{m-1}{2}]$,


$[-1+i,1-i]$ и боковыми криволинейными сторонами ${L}_m$ и


$\widetilde{L}_m$, когда $0<m<1$. Уравнение линии


${L}_m:\;\beta=\frac{(1+m)^2\alpha}{4m\alpha-(1-m)^2}$ $($соединяет


точки $\frac{m-1}{2}+i\frac{m+1}{2}$ и $-1+i)$, уравнение линии


$\widetilde{L}_m:\;\beta=\frac{(1-m)^2\alpha}{4m\alpha-(1+m)^2}$


$($соединяет точки $\frac{m+1}{2}+i\frac{m-1}{2}$ и $1-i)$.


Расширить эти области нельзя, т.\,к. знаки равенства в $(\ref{23})$


осуществляются на функциях с


$f'_0(\zeta)=\frac{1+\beta\zeta^2}{1-\alpha\zeta^2}$.


}


\medskip





С помощью этой леммы будет доказана





\begin{thm}


Если регулярная в $E$ функция $f(\zeta)$ имеет


производную, удовлетворяющую подчинению


$$


f'(\zeta)\prec T(\alpha,\beta;\zeta),\ \; \alpha+i\beta\in {\cal


D}_m\quad ({\cal D}_m \text{ --- из текста леммы}),


$$


и $f''(0)=0$, то в случае функции ${\cal F}(\zeta)$ со связью


${\cal F}'(\zeta)=f'{}^{c}(\zeta)$, содержащей комплексную


постоянную $c$, конформный радиус области ${\cal F}(E)$ будет


иметь единственную критическую точку в $0$ при условии


\begin{equation}\label{24}


m|c|<1.


\end{equation}


\end{thm}





{\bf Доказательство.} Так как ${\cal


F}''(0)=cf'{}^{c-1}(0)f''(0)=0$ , то уравнение (\ref{10}),


записанное для ${\cal F}(\zeta)$, будет удовлетворено при


$a=0$. Чтобы доказать отсутствие других критических точек у


конформного радиуса области ${\cal F}(E)$, выведем неравенство


\begin{equation}\label{25}


\bigg|\frac{{\cal F}''(\zeta)}{{\cal


F}'(\zeta)}\bigg|< \frac{2|\zeta|}{1-|\zeta|^2}\ \;


\text{ при \;} \zeta\neq 0.


\end{equation}





Для этого воспользуемся логарифмической производной соотношения ${\cal


F}'(\zeta)=[f'(\zeta)]^c$ в форме


$$


\frac{{\cal F}''}{{\cal F}'}= c\frac{f''}{f'}.


$$


Подставив это равенство в (\ref{23}), с учетом (\ref{24}) получим


$$


\bigg|\frac{{\cal F}''(\zeta)}{{\cal


F}'(\zeta)}\bigg|\leq \frac{2m|c|\,|\zeta|}{1-|\zeta|^2}\ \Rightarrow


\ \text{(\ref{25})}.\quad\square


$$





Определяющую роль в теореме 3 играют подчиняющие функции 


$T(\alpha,\beta;\zeta)=(1+\beta\zeta)/(1-\alpha\zeta)$, которые легко 


упорядочить следующим образом: 


\begin{multline} 


T(\alpha,\beta;\zeta)\prec T(\alpha_0,\beta_0;\zeta) \Leftrightarrow 


(T(\alpha,\beta;1)\le T(\alpha_0,\beta_0;1), 


T(\alpha,\beta;-1)\ge T(\alpha_0,\beta_0;-1)) \Leftrightarrow \\ 


\Leftrightarrow\{\alpha(1+\beta_0)+\beta(1-\alpha_0)\le 


\alpha_0+\beta_0,\ \alpha(1-\beta_0)+\beta(1+\alpha_0)\le 


\alpha_0+\beta_0\}.


\end{multline} 


Неравенства (27) отражают наглядный геометрический факт, который можно 


пояснить с помощью диффеоморфного соответствия ([7], с.\,47) 


\begin{equation} 


a=\frac{1+\alpha\beta}{1-\alpha^2},\ \; 


b=\frac{\alpha+\beta}{1-\alpha^2}\; \Leftrightarrow\; 


\alpha=\frac{a-1}{b},\ \; \beta=b-\frac{a(a-1)}{b}. 


\end{equation} 


Именно, критерием того, что круг $E(a_0,b_0)=\{w:|w-a_0|<b_0\}$ с 


вещественным центром $a_0$ содержит круг $E(a,b)$ с вещественным 


центром $a$, являются два неравенства 


\begin{equation} 


a+b\le a_0+b_0,\ \; a-b\ge a_0-b_0\ \; (\Leftrightarrow\; 


E(a,b)\subseteq E(a_0,b_0)), 


\end{equation} 


в которых задействованы точки пересечения окружностей $\partial E(a,b)$ 


и $\partial E(a_0,b_0)$ с вещественной осью. 


 


Неравенства (29) определяют две полуплоскости в плоскости параметров 


$a+ib$. Аналогичные полуплоскости, которые обозначим через 


$\Pi_1(\alpha_0,\beta_0)$ и $\Pi_2(\alpha_0,\beta_0)$, в плоскости 


$\alpha+i\beta$ являются изображениями неравенств (27), которые можно 


получить из (29) левой подстановкой в (28). Граничная прямая 


полуплоскости $\Pi_1(\alpha_0,\beta_0)$ проходит через точки 


$\alpha_0+i\beta_0$ и $1-i$, а прямая $\partial\Pi_2(\alpha_0,\beta_0)$ 


задается точками $\alpha_0+i\beta_0$ и $-1+i$. Обеим полуплоскостям 


принадлежит начало координат. Пересечение $\Pi_1(\alpha_0,\beta_0)\cap 


\Pi_2(\alpha_0,\beta_0)\cap \cal D_m$ определяет область значений параметров 


в плоскости $\alpha+i\beta$, для которых теорема 3 дает результат 


слабее, чем при $\alpha_0+i\beta_0$. Поэтому наилучшие эффекты связаны 


с условиями подчинения в (24) и в теореме 3, когда $m<1/|c|$ и


\begin{equation} 


\alpha+i\beta\in\partial \cal D_m\setminus[-1+i,1-i], 


\end{equation} 


т.\,е. $\alpha+i\beta$ лежит на границе области $\cal D_m$ с исключением 


отрезка, соединяющего точки $-1+i$ и $1-i$. Следовательно, в первой 


части импликации (24) вместо $\alpha+i\beta\in \cal D_m$ можно вставить 


(30). 


 


Заметим, что в [4] была обоснована единственность критической точки 


конформного радиуса области $f(E)$ при условии


$$


\{f'{}^{\Tilde a}(\zeta)\prec 


T(\alpha,\beta;\zeta),\ \alpha+i\beta\in \cal D_m\}\ \;


\text{и} \ \; m\le|\wt a|.


$$


Здесь также включение $\alpha+i\beta\in \cal D_m$ можно заменить на (30). 


Аналогичные улучшения можно внести с помощью включения вида (30) для 


области $\cal G_m$ в теорему 2 из [4]. 


 


Другие виды условий, гарантирующих единственность критической точки 


конформного радиуса, и приложения этих условий к внешним обратным 


краевым задачам [10], [11] будут изложены в другой статье авторов. 


 





\begin{thebibliography}{99}





\bibitem{1}


Аксентьев~Л.А. {\em Связь внешней обратной краевой задачи с внутренним


радиусом области} // Изв. вузов. Математика. -- 1984. -- \No\,2.


-- С.\,3--11.





\bibitem{2}


Krzy\v{z} J.G. {\em Some remarks on the maxima of inner conformal


radius} // Ann. UMCS. A. -- 1992. -- V.\,46. -- P.\,57--61.





\bibitem{3}


K\"uhnau R. {\em Maxima beim konformen Radius einfach


zusammenh\"angender Gebiete} // Ann. UMCS. A. -- 1992. - V.\,46. --


P.\,63--73.





\bibitem{4}


Аксентьев Л.А., Казанцев А.В., Попов Н.И. {\em О теоремах


единственности для внешней обратной краевой задачи в подклассах


однолистных функций} // Изв. вузов. Математика. -- 1998. -- \No\,8.


-- С.\,3--13.





\bibitem{5}


Аксентьев Л.А., Казанцев А.В. {\em Новое свойство класса Нехари и его


применение} // Изв. вузов. Математика. -- 1989. -- \No\,8.


-- С.\,69--72.





\bibitem{6}


Nehari Z. {\em The Schwarzian derivative and schlicht functions} // Bull.


Amer Math. Soc. -- 1949. -- V.\,55. -- \No\,6. -- P.\,545--551.





\bibitem{7}


Аксентьев Л.А., Казанцев А.В., Киндер М.И., Киселев А.В. {\em О классах


единственности внешней обратной краевой задачи} // Тр. семин. по


краев. задачам. -- Казань: Изд-во Ка\-занск. ун-та, 1990. -- Вып.\,24.


-- С.\,39--62.





\bibitem{8}


Голузин Г.М. {\em Геометрическая теория функций комплексного


переменного}. -- 2-е изд. -- М.: Наука, 1966. -- 628 с.





\bibitem{9}


Okuyama Y\^{u}suke. {\em The norm estimates of pre-Schwarzian


derivatives of spiral-like functions} / International conference on


complex analysis and related topics (the VIII-th romanian-finnish


seminar). -- August 23--27, 1999. -- Jassy, Romania. -- 1999. -- P.\,52--53.





\bibitem{10}


Гахов Ф.Д. {\em Краевые задачи}. -- 3-е изд. -- М.: Наука, 1977. -- 640 с.





\bibitem{11}


Тумашев Г.Г., Нужин М.Т. {\em Обратные краевые задачи и их приложения}.


-- 2-е изд. -- Казань: Изд-во Казанск. ун-та, 1965. -- 333 с.





\end{thebibliography}





\vskip 2cm





\post{\it Казанский государственный университет }{\it Поступила}


\post{\it Марийский государственный университет }{07.02.2000}





\label{end}


\end{document}





